\chapter{Complete CRUD: The Task Resource End to End}

This chapter wires together everything from Chapters 9-13 into a full, working CRUD
feature for the Task resource — model, routes, controller, service, and the React
pieces that consume the API.

\section{Backend: Routes}

\begin{lstlisting}[language=JavaScript]
// routes/task.routes.js
const express = require("express");
const {
  createTask,
  getTasks,
  getTaskById,
  updateTask,
  deleteTask,
} = require("../controllers/task.controller");
const { protect } = require("../middleware/auth.middleware");

const router = express.Router();

router.use(protect);

router.post("/", createTask);
router.get("/", getTasks);
router.get("/:id", getTaskById);
router.put("/:id", updateTask);
router.delete("/:id", deleteTask);

module.exports = router;
\end{lstlisting}

\section{Backend: Service Layer}

\begin{lstlisting}[language=JavaScript]
// services/task.service.js
const Task = require("../models/Task");

const createTask = async (data, userId) => {
  return Task.create({ ...data, user: userId });
};

const getTasks = async (userId, { page = 1, limit = 10, status }) => {
  const filter = { user: userId };
  if (status) filter.status = status;

  const tasks = await Task.find(filter)
    .sort({ createdAt: -1 })
    .skip((Number(page) - 1) * Number(limit))
    .limit(Number(limit));

  const total = await Task.countDocuments(filter);
  return { tasks, total, page: Number(page), limit: Number(limit) };
};

const getTaskById = async (id, userId) => {
  const task = await Task.findOne({ _id: id, user: userId });
  if (!task) {
    const err = new Error("Task not found");
    err.statusCode = 404;
    throw err;
  }
  return task;
};

const updateTask = async (id, userId, data) => {
  const task = await Task.findOneAndUpdate(
    { _id: id, user: userId },
    data,
    { new: true, runValidators: true }
  );
  if (!task) {
    const err = new Error("Task not found");
    err.statusCode = 404;
    throw err;
  }
  return task;
};

const deleteTask = async (id, userId) => {
  const task = await Task.findOneAndDelete({ _id: id, user: userId });
  if (!task) {
    const err = new Error("Task not found");
    err.statusCode = 404;
    throw err;
  }
  return task;
};

module.exports = { createTask, getTasks, getTaskById, updateTask, deleteTask };
\end{lstlisting}

\section{Backend: Controller}

\begin{lstlisting}[language=JavaScript]
// controllers/task.controller.js
const taskService = require("../services/task.service");

const createTask = async (req, res, next) => {
  try {
    const task = await taskService.createTask(req.body, req.user.id);
    res.status(201).json(task);
  } catch (err) {
    next(err);
  }
};

const getTasks = async (req, res, next) => {
  try {
    const result = await taskService.getTasks(req.user.id, req.query);
    res.status(200).json(result);
  } catch (err) {
    next(err);
  }
};

const getTaskById = async (req, res, next) => {
  try {
    const task = await taskService.getTaskById(req.params.id, req.user.id);
    res.status(200).json(task);
  } catch (err) {
    next(err);
  }
};

const updateTask = async (req, res, next) => {
  try {
    const task = await taskService.updateTask(req.params.id, req.user.id, req.body);
    res.status(200).json(task);
  } catch (err) {
    next(err);
  }
};

const deleteTask = async (req, res, next) => {
  try {
    await taskService.deleteTask(req.params.id, req.user.id);
    res.status(200).json({ message: "Task deleted" });
  } catch (err) {
    next(err);
  }
};

module.exports = { createTask, getTasks, getTaskById, updateTask, deleteTask };
\end{lstlisting}

\section{Frontend: taskService.js}

\begin{lstlisting}[language=JavaScript]
// src/services/taskService.js
import api from "./api"; // preconfigured axios instance with baseURL + credentials

export const fetchTasks = (params) => api.get("/tasks", { params });
export const fetchTaskById = (id) => api.get(`/tasks/${id}`);
export const createTask = (data) => api.post("/tasks", data);
export const updateTask = (id, data) => api.put(`/tasks/${id}`, data);
export const deleteTask = (id) => api.delete(`/tasks/${id}`);
\end{lstlisting}

\section{Frontend: TaskForm.jsx}

\begin{lstlisting}[language=JavaScript]
import { useState } from "react";
import { createTask, updateTask } from "../services/taskService";

const initialState = { title: "", description: "", status: "todo", priority: "medium", dueDate: "" };

const TaskForm = ({ existingTask, onSaved }) => {
  const [form, setForm] = useState(existingTask || initialState);
  const [saving, setSaving] = useState(false);

  const handleChange = (e) => {
    setForm({ ...form, [e.target.name]: e.target.value });
  };

  const handleSubmit = async (e) => {
    e.preventDefault();
    setSaving(true);
    try {
      if (existingTask) {
        await updateTask(existingTask._id, form);
      } else {
        await createTask(form);
      }
      onSaved();
    } finally {
      setSaving(false);
    }
  };

  return (
    <form onSubmit={handleSubmit}>
      <input name="title" value={form.title} onChange={handleChange} placeholder="Title" required />
      <textarea name="description" value={form.description} onChange={handleChange} placeholder="Description" />
      <select name="status" value={form.status} onChange={handleChange}>
        <option value="todo">To Do</option>
        <option value="progress">In Progress</option>
        <option value="completed">Completed</option>
      </select>
      <select name="priority" value={form.priority} onChange={handleChange}>
        <option value="low">Low</option>
        <option value="medium">Medium</option>
        <option value="high">High</option>
      </select>
      <input type="date" name="dueDate" value={form.dueDate} onChange={handleChange} />
      <button type="submit" disabled={saving}>{saving ? "Saving..." : "Save Task"}</button>
    </form>
  );
};

export default TaskForm;
\end{lstlisting}

\section{Frontend: TaskList.jsx}

\begin{lstlisting}[language=JavaScript]
import { useEffect, useState } from "react";
import { fetchTasks, deleteTask } from "../services/taskService";
import TaskCard from "./TaskCard";

const TaskList = () => {
  const [tasks, setTasks] = useState([]);
  const [loading, setLoading] = useState(true);

  const loadTasks = async () => {
    setLoading(true);
    const { data } = await fetchTasks({ page: 1, limit: 10 });
    setTasks(data.tasks);
    setLoading(false);
  };

  useEffect(() => {
    loadTasks();
  }, []);

  const handleDelete = async (id) => {
    await deleteTask(id);
    setTasks((prev) => prev.filter((t) => t._id !== id));
  };

  if (loading) return <p>Loading tasks...</p>;

  return (
    <div className="task-list">
      {tasks.map((task) => (
        <TaskCard key={task._id} task={task} onDelete={() => handleDelete(task._id)} />
      ))}
    </div>
  );
};

export default TaskList;
\end{lstlisting}

\section{End-to-End Flows}

\begin{diagrambox}[Create Task Flow]
\begin{verbatim}
TaskForm submit
     |
     v
createTask(form)  ------> POST /api/tasks  (axios, withCredentials)
     |
     v
protect middleware -----> verifies auth cookie, sets req.user
     |
     v
task.controller.createTask -> reads req.body, calls service
     |
     v
task.service.createTask ---> Task.create({...data, user: userId})
     |
     v
MongoDB inserts document
     |
     v
201 response -----------> onSaved() -> TaskList reloads
\end{verbatim}
\end{diagrambox}

\begin{diagrambox}[Delete Task Flow]
\begin{verbatim}
TaskCard "Delete" click
     |
     v
deleteTask(id)  --------> DELETE /api/tasks/:id
     |
     v
protect middleware -----> req.user set
     |
     v
task.controller.deleteTask -> calls service with id + userId
     |
     v
task.service.deleteTask --> Task.findOneAndDelete({_id, user})
     |
     v
200 response ------------> setTasks(prev => filter out id) (optimistic UI update)
\end{verbatim}
\end{diagrambox}

\begin{warnbox}[WARNING]
Every service function scopes queries with \texttt{user: userId}. Without this, any authenticated user could read, update, or delete another user's tasks just by guessing an \texttt{\_id}.
\end{warnbox}
